% Conclusion (draft).
\section{Conclusion}
\label{sec:conclusion}

We formulated the offloading of DAG workflows in cloud--edge--end systems with interval-valued resources and workloads
as an interval multi-objective problem over task placement and task order, with makespan computed along the interval
critical path, device energy and cost as objectives. Extending the three-tier selection of IMOEA-DQN to this problem
showed that its convergence ranking, the interval dominance degree, stops selecting once most of the population is
mutually non-dominated, and that its learned control of the selection weights does not beat a fixed choice. \ours{}
replaces the convergence ranking with an interval version of the MSOPS best-position Tchebycheff rank, breaks its ties
by crowding distance, mates neighbours in objective space and fixes the selection weights on training instances.

On Jin et al.'s instance sizes \ours{} has the best IGD and hypervolume on every instance, and on five families of
scientific workflows with 100 to 1000 tasks it is better than the strongest baseline, IMOEA/D, on 43 of 45 instances in
IGD and 42 in hypervolume, at about the same running time. The ablation attributes the gain to the new rank and the
neighbourhood mating together, and the predicted objective intervals contain every realised value in a Monte Carlo
check. At 1000 and 2000 tasks the advantage narrows: \ours{} keeps the higher hypervolume, but IMOEA/D reaches cheaper
solutions at the cost end of the front and IGD becomes level and dependent on the reference front.

Three directions follow. The first is to close the gap at the cost extreme on large workflows, for example by
combining the three-tier selection with steady-state replacement per subproblem, which our analysis points to as the
source of IMOEA/D's advantage there. The second is to make the model more realistic: shared link bandwidth, receive
energy, time-varying edge capacity, correlated uncertainty and applications that arrive online. The third is a small
uncertainty share in the selection weights, which narrowed the intervals by 10\% at a small cost in IGD on the training
instances and may suit settings where predictable performance matters more than the best expected value.
